\chapter{Bruises (Contusions and Ecchymosis)}

\section*{Definition}
A bruise (contusion or ecchymosis) is a discoloration of the skin caused by bleeding into the subcutaneous tissue from ruptured blood vessels following blunt trauma. Bruises evolve through characteristic color changes as hemoglobin is broken down and reabsorbed.

\section*{When to See a Doctor}
\begin{itemize}
\item Large or painful bruises without significant trauma
\item Frequent or recurrent bruising with minimal provocation
\item Bruises accompanied by prolonged bleeding from cuts or gums
\item Spontaneous bruising (no identifiable cause)
\item Bruising with nosebleeds, blood in urine or stool
\item Petechiae (tiny red spots indicating platelet disorders)
\item Bruising in unusual locations
\item Bruising in a child with an inconsistent explanation (possible abuse)
\end{itemize}

\section*{How It Develops}
Bruising occurs when damaged small blood vessels leak blood into nearby tissue. The color and size may change as the blood is broken down and reabsorbed, but the timing and color sequence are variable. Most uncomplicated bruises fade over about 1--3 weeks.

\section*{Causes}
\begin{itemize}
\item Accidental trauma (most common cause)
\item Anticoagulant therapy (warfarin, heparin, DOACs)
\item Antiplatelet therapy (aspirin, clopidogrel)
\item Coagulation disorders (hemophilia, von Willebrand disease)
\item Thrombocytopenia (low platelet count)
\item Corticosteroid use (skin thinning and vascular fragility)
\item Vitamin C deficiency (scurvy)
\item Vitamin K deficiency
\item Liver disease
\item Aging (increased vascular fragility)
\item Physical abuse
\end{itemize}

\section*{Related Symptoms}
\begin{itemize}
\item Petechiae (tiny capillary hemorrhages)
\item Hematoma (larger collection of blood in tissue)
\item Purpura (bleeding into skin)
\item nosebleed (nosebleeds)
\item Gingival bleeding (bleeding gums)
\item Hemarthrosis (bleeding into joints)
\item Coagulopathy (bleeding tendency)
\item Thrombocytopenia
\end{itemize}

\section*{Medical Evaluation}
\begin{itemize}
 \item Complete blood count, including platelet count, when a bleeding disorder is suspected
 \item PT/INR and aPTT when indicated by the bleeding history, medicines, or examination
 \item Further tests for von Willebrand disease or clotting-factor disorders if the initial assessment suggests them
 \item Liver or other targeted tests when the history or examination suggests an underlying disease
 \item Imaging if a deep hematoma, significant injury, or fracture is suspected
\end{itemize}

\section*{Treatment Options}
\begin{itemize}
 \item Treat the underlying cause; reversal or factor treatment is reserved for selected serious bleeding and is clinician-directed
 \item Do not stop anticoagulant or antiplatelet medicines without medical advice; the clinician weighs bleeding against clotting risk
\item Platelet transfusion for severe thrombocytopenia
\item Surgical evacuation for large, tense hematomas
\item Counseling and social services if abuse is suspected
\end{itemize}

\section*{Self-Care and Home Management}
\begin{itemize}
\item Apply ice wrapped in a cloth for 15-20 minutes to reduce swelling
\item Elevate the injured area if possible
\item Rest the injured area to prevent further trauma
\item Apply a compression bandage to minimize swelling
\item Avoid heat for the first 48 hours
 \item Evidence for arnica and bromelain is limited; do not apply products to broken skin
 \item Ask a clinician or pharmacist before using aspirin or NSAIDs if bruising is unexplained or frequent
\end{itemize}

\section*{References}
\begin{thebibliography}{99}
\bibitem{bruises:ref1} Liem RI, Lanzkron S, et al. ``Approach to the patient with bruising.'' In: Hematology: Basic Principles and Practice. 7th ed. Elsevier; 2018.
\bibitem{bruises:ref2} Neff AT. ``Autoimmune and acquired coagulation disorders.'' Hematology Am Soc Hematol Educ Program. 2020;2020(1):546-554.
\bibitem{bruises:ref3} Rodeghiero F, Pabinger I, Ragni M, et al. ``Fundamentals for a systematic approach to mild bleeding disorders.'' Haemophilia. 2012;18(Suppl 4):64-69.
\bibitem{bruises:ref4} James AH, Manco-Johnson MJ, Yawn BP, et al. ``Von Willebrand disease: key points from the 2021 NIH/CDC guidelines.'' J Womens Health. 2022;31(2):179-186.
\bibitem{bruises:ref5} Kessler CM. ``Diagnosis and management of inherited bleeding disorders.'' UpToDate. Wolters Kluwer; 2025.
\bibitem{bruises:ref6} Kumar V, Abbas AK, Aster JC. ``Hemorrhagic disorders.'' In: Robbins and Cotran Pathologic Basis of Disease. 10th ed. Elsevier; 2021.
\end{thebibliography}
